Dane:

- Typy statków:

\(typy \in {1,2,3}\)

- Liczba statków:
\(S = 12\)

- Liczba doków:

\(D = 2\)

- Liczba miejsc do cumowania w dokach:

\(C=10\)

- Lista statków:
\(L=[...]\)

Zmienne decyzyjne:

\begin{equation}
    M[N][D][C]
\end{equation}


Funkcja celu:
%  bez premiowania statków
% \begin{equation}
%     max \sum_{s=1}^{S} \sum_{d=1}^{D} \sum_{c=0}^{C} M[s][d][c]
% \end{equation}

\begin{equation}
    max \sum_{s=1}^{S} L[s] * (\sum_{d=1}^{D} \sum_{c=0}^{C} M[s][d][c])
\end{equation}

% MOŻLIWĄ koncepkcją jest też to aby zrobić dwie funkcje celu, jedna z premiowaniem statków druga bez premiowania. Wtedy przy maksymalizacji obu będziemy mieć po pierwsze najwięcej wartościowych statków, a w drugiej kolejności po prostu najwięcej statków w porcie.

Ograniczenia:

\begin{equation} \label{eq:uniq}
    \forall_{s\in S} ( \forall_{d\in D}\forall_{c \in C} M[n][d][s] \leq 1 )
\end{equation}

W \ref{eq:uniq} zapewniamy, że statek, będzie co najwyżej raz w porcie. 

\begin{equation} \label{eq:no-duplication}
    \forall_{d\in D}\forall_{c \in C} (\forall_{s \in S} M[n][d][s] \leq 1)
\end{equation}

W \ref{eq:no-duplication} zapewniamy, że nie ma dwóch statków w tym samym miejscu.

\begin{equation} \label{eq:3jka-blokuje}
    \begin{split}
        ((\forall_{s_1\in S}\forall_{d_1\in D}\forall_{c_1 \in C} M[s_1][d_1][c_1] \equiv 1) \land (L[s_1] \equiv 3 )) \implies  \\
    (\forall_{s_2>s_1} \forall_{d_2 \in D} \forall_{c_2 \leq c} M[s_2][d_2][c_2] \equiv 0)
    \end{split}
\end{equation}

W \ref{eq:3jka-blokuje} zapewniamy, że statek o wielkości 3 blokuje część wewnętrzną portu.

\begin{equation} \label{eq:2jka-blokuje}
    \begin{split}
        (((\forall_{s_1\in S}\forall_{d_1\in D}\forall_{c_1 \in C}  M[s_1][d_1][c_1] \equiv 1) \land (L[s_1] = 2 )) \land \\
        ( \exists_{s_2} M[s_2][ \neg d_1][c_1] = 1) \land ( L[s_2] \equiv 2 ))  ) \\
        \implies (\forall_{s_3>s_1} \forall_{d_2 \in D} \forall_{c_3 \leq c} M[s_3][d][c_3] \equiv 0)
    \end{split}
\end{equation}

W \ref{eq:2jka-blokuje} zapewniamy sobie, że dwa statki na przeciwko siebie również blokują wewnętrzną stronę portu.

